% The goals as Lean statements, design decisions, stretch goals, a summary
% and further reading
\begin{frame*}
  \frametitle{Each Level Ends in a Theorem}
  {\fvset{fontsize = \footnotesize}\leanblock{goals.lean}}
  \sketch
  \worktag{Levels 3--5}
\end{frame*}

\begin{frame}
  \frametitle{Your Design Decides How Hard the Proofs Are}
  \begin{itemize}
    \item \textbf{Attributes}: named or positional? Names make
      $R \times S = S \times R$ true.
    \item \textbf{Values}: one type for all columns, or a type per
      column?
    \item \textbf{Bad queries}: rejected when evaluated, or impossible to
      write?
    \item \textbf{Relations}: lists, bags or sets?
    \item \textbf{Conditions}: syntax you can inspect, for the side
      conditions.
  \end{itemize}
  \begin{sticky}[halign = center]
    No design is right first time: if the proofs hurt, change the
    definitions.
  \end{sticky}
  \slideref{\lnk{fpil-universe} \rfsep \lnk{fpil-pitfalls}}
\end{frame}

\begin{frame}
  \frametitle{Beyond Level 5: Bags, NULL, Aggregation, SQL}
  \begin{itemize}
    \item \textbf{Bags}: SQL keeps duplicates. Which laws survive?
    \item \textbf{NULL}: missing values, and three truth values.
    \item \textbf{Aggregation}: \code{GROUP BY}, \code{COUNT},
      \code{SUM}.
    \item \textbf{SQL}: translate a fragment of SQL into your algebra.
    \item \textbf{Codd's theorem}: the algebra and (safe) first-order
      logic express the same queries.
    \item \textbf{Containment}: decide whether a conjunctive query's
      answers are always among another's.
  \end{itemize}
  \slideref{\lnk{fpil-tqex} \rfsep
    A. K. Chandra, P. M. Merlin. \textbf{Optimal Implementation of Conjunctive Queries in Relational Data Bases}. STOC 1977.}
  \stretchtag
\end{frame}

\begin{frame}
  \frametitle{The Lean You Need, at a Glance}
  {\centering\renewcommand\arraystretch{1.35}
    {\headerrow\begin{tabular}{@{}l@{\qquad}>{\raggedright\arraybackslash}p{62mm}@{\qquad}>{\raggedright\arraybackslash}p{46mm}@{}}
      \hline
      \textbf{Level} & \textbf{Lean} & \textbf{Read first} \\ \hline
      1 Tables & structures, inductive types, lists & FPIL 1, 3.5 \\
      2 Queries & recursion; indexed families & FPIL 7.1--7.3; TPIL 7, 8 \\
      3 Laws & propositions, tactics, sets & FPIL interludes; TPIL 3--5 \\
      4 Optimiser & induction on your own types & TPIL 7.6, 8 \\
      5 Fast joins & arrays, hash maps, equivalence & FPIL 8 \\ \hline
    \end{tabular}}\par}
  \slideref{\lnk{fpil} \rfsep \lnk{tpil} \rfsep \lnk{mil}}
\end{frame}

\begin{frame}
  \frametitle{Others Did This in Coq: Read Them After Trying}
  Verified relational databases, in Coq (now Rocq):
  \begin{itemize}
    \item G. Malecha, G. Morrisett, A. Shinnar, R. Wisnesky. \emph{Toward a
      Verified Relational Database Management System}. POPL 2010.
    \item V. Benzaken, \'E. Contejean, S. Dumbrava. \emph{A Coq
      Formalization of the Relational Data Model}. ESOP 2014.
    \item S. Chu, K. Weitz, A. Cheung, D. Suciu. \emph{HoTTSQL: Proving
      Query Rewrites with Univalent SQL Semantics}. PLDI 2017.
    \item V. Benzaken, \'E. Contejean. \emph{A Coq Mechanised Formal
      Semantics for Realistic SQL Queries}. CPP 2019.
  \end{itemize}
  For the theory of databases: \lnk{ahv}. Addison-Wesley, 1995.
\end{frame}
